% Transcription — fonds Alexandre Grothendieck, shelfmark 162-5, batch 2 (pages 21-40).
% Established from the facsimile archives/batches/162-5/batch-02.pdf (a mirror of
% https://grothendieck.umontpellier.fr/162-5.pdf), per the skill
% transcribe-grothendieck. The batch file carries no cover sheet: PDF page N is
% archivists' page 20 + N, checked against the pencilled numbers bottom-left.
%
% Pass: Opus 5.5 (claude-opus-5-5), 2026-09-22 — first pass, unchecked against the pages by a human.
%
% What the batch holds, in the archivists' order.
%
%   21-30  Ten leaves of manuscript in blue ink, headed top right in his hand
%          « tapis Quillen (Notes du 6.9.68) » — the « notes manuscrites
%          (1968) » of the inventory. Numbered by paragraphs 1 to 11, no
%          pagination of his own. Formal categories over a scheme X_0 (a
%          pro-ring A with augmentation and a diagonal Δ), the conormal Ω and
%          the De Rham differential d; general De Rham complexes (a graded
%          differential on ΛΩ, determined by δ_0, δ_1 under α), β), γ));
%          the Lie ring g = Hom(Ω, O) with θ_X and the bracket (*), (**)
%          θ_[X,Y] = [θ_X, θ_Y], and a Jacobi computation left open
%          (« pourquoi est-ce nul ?? », p. 25); the cochain complex C*(g);
%          Quillen's theorem (char. 0): formal categories with Ω flat ≃ De
%          Rham complexes with Ω flat, containing Cartan's theorem on formal
%          groups, with a divided-power variant; coefficients: descent data
%          ↔ connections δ_M, curvature, integrable g-connections, and the
%          quasi-isomorphism ΛΩ ⊗ M ≃ C*(C, M) between De Rham and
%          Hochschild (Čech-Alexander) cohomology. The last line, « Variante
%          à puissances divisées », opens a paragraph that is not written.
%   31     Blank separator sheet: absent.
%   32-40  A second copy of the typescript « Tapis de Quillen », dated in his
%          hand « notes du 10.9.68 », its typed pages 1 to 9 — the text
%          batch 1 gives from the other copy at pages 2-10: part I entire and
%          part II up to K_n(C) and the tapis de Roos. The copy runs on into
%          batch 3 (page 41 finishes the sentence broken off at the foot of
%          page 40, then part III). Its ink layer is its own: a framed note
%          heading page 38 on Hom_grab and RHom_Z(C_x, C'_x), the ink
%          insertion and the struck commutative variants on page 38 with the
%          margin explaining why (a commutative Picard category is realised
%          only by a non-commutative (Cat)-group), the answer « non … pas
%          essentiellement surjectif » to the loop-space question on page 39,
%          the « Attention » on étale topologies at the foot of page 34.
%          Following batch 3 for this same copy, typed text including typed
%          interlinear insertions is given plain; ink insertions are \add,
%          words struck in ink \struck, detached ink notes \marginal;
%          typewriter x-outs are not reproduced; typed underlining is \emph.
%          His hand turns the typed superscripts of C_0, C_1, X_ij into
%          subscripts; the correction is followed (as in batch 1).
%
% Nothing in this batch is a printed offprint; no page was skipped as
% unrelated. Only page 31 is absent.
%
% Notation fixed for the batch. In the manuscript: his gothic g is
% \mathfrak{g}; the exterior algebra with a star over Λ is
% \overset{*}{\Lambda}\Omega, its degree-p part \overset{p}{\Lambda}; O_{X_0}
% is \mathcal{O}_{X_0} (he writes O_X and O_{X_0} indifferently, left as
% written); underlined Hom is \underline{\mathrm{Hom}}; θ_X, i_X, δ as
% written, δ_M^0, δ_M^1 for the components of a connection; C for the formal
% category is \mathcal{C}, the cochain complexes C^p, C^*(g) plain. In the
% typescript, as batch 1: C° is C^\circ, « lim » over a left arrow is
% \varprojlim, the tilde topos \widetilde{C}, D^b_lc as D^b_{\mathrm{lc}},
% Fl(C) underlined \underline{\mathrm{Fl}}(C), the symmetric group Sym_n.
%
% Slips left as written, flagged where they bear on the mathematics: β)
% p. 22 (δ_0(f)δ_1(ω) for δ_0(f)∧ω), « section X∧Y » p. 23, the missing
% sign in the Leibniz rule p. 28, K'(u)(X,X) p. 29 (his own margin says
% « Attention signe ! »). Typed slips (« coéfficients », « exctitude »,
% « catétorie », « satsifont », « équivelente », « homomorhisme »,
% « symmétrisé », « flêches », « il f et s ») are left as typed.
%
% Doubtful, not settled: much of the connective prose of pages 21-22 and
% 26-27; the struck margin block and the first half of the margin note on
% page 38; the phrase between the formulas of page 29; the last line of the
% slanted margin of page 30; « une équivalence » at the head of page 35,
% whose typed context is overstruck.
%
% For the modernised reading. The manuscript (6.9.68) and the typescript
% (10.9.68) are two stages four days apart: section 11 of the manuscript
% (De Rham vs Hochschild cohomology, « Relier au th. fondamental du
% dévissage des cristaux … topos convenable ?! ») is the crystalline side
% the typescript does not take up. The Jacobi identity of page 25 is left
% as a question by him (« pourquoi est-ce nul ?? »); it is not settled on
% the page. The dating is the inventory's own, for the shelfmark.
\documentclass[11pt,a4paper]{article}
% Le préambule commun à toutes les transcriptions du fonds Grothendieck.
%
% Il tient en une idée : **l'incertitude se compose, elle ne se résout pas.**
% Une transcription qui lisse un mot douteux a détruit la seule chose pour
% laquelle on la faisait — un lecteur doit pouvoir savoir, à chaque endroit,
% ce qui était sur le feuillet et ce qui a été ajouté. Les six macros qui
% suivent sont donc l'appareil critique, et non de la décoration.
%
% Les mêmes commandes sont reconnues par `scripts/render.mjs`, qui produit la
% vue de lecture du site. Une macro ajoutée ici doit l'être aussi là-bas,
% faute de quoi le rendu échouera bruyamment — ce qui est voulu : un rendu
% silencieusement faux est pire qu'un rendu absent.

\ProvidesPackage{grothendieck}[2026/08/08 Transcriptions du fonds Grothendieck]

\usepackage{amsmath,amssymb,amsthm}
\usepackage{mathrsfs}
\usepackage{stmaryrd}
% Les diagrammes commutatifs sont partout dans ce fonds : c'est la langue
% même de la géométrie algébrique de Grothendieck, et une transcription qui
% les rendrait en prose ne transcrirait rien.
\usepackage{tikz-cd}
% Français par défaut — c'est la langue des feuillets — mais l'anglais est
% chargé aussi : la traduction et le résumé sont des documents anglais, et la
% typographie française (espaces avant les deux-points) y serait une faute.
% Ces documents basculent par \selectlanguage{english} en tête de corps.
\usepackage[english,french]{babel}
\usepackage{fontspec}
\usepackage{geometry}
\geometry{a4paper,margin=25mm,marginparwidth=28mm}
\usepackage{marginnote}
\usepackage{xcolor}
% ulem, not soul. `soul`'s \st and \ul are built on a character-by-character
% scan that XeLaTeX's font machinery breaks: the document fails with an
% undefined control sequence rather than degrading. ulem does the same two jobs
% and is Unicode-engine safe. `normalem` stops it hijacking \emph, which the
% transcriptions use for Grothendieck's own emphasis.
\usepackage[normalem]{ulem}
\usepackage{hyperref}
\hypersetup{colorlinks=true,linkcolor=black,urlcolor=black,pdfborder={0 0 0}}

% --- Métadonnées du lot -----------------------------------------------------
% Reprises telles quelles par le script de rendu, qui les affiche en tête de
% la vue de lecture. Elles fixent ce que le fichier prétend couvrir : sans
% elles, une transcription flotte sans point d'attache dans un fonds de seize
% mille feuillets.

\newcommand{\folder}[1]{\gdef\grothFolder{#1}}
\newcommand{\batch}[1]{\gdef\grothBatch{#1}}
\newcommand{\leaves}[2]{\gdef\grothFirst{#1}\gdef\grothLast{#2}}
\newcommand{\foldertitle}[1]{\gdef\grothFoldertitle{#1}}
\gdef\grothFolder{?}\gdef\grothBatch{?}\gdef\grothFirst{?}\gdef\grothLast{?}\gdef\grothFoldertitle{}

% --- Le résumé --------------------------------------------------------------
%
% Il ouvre la lecture modernisée, sous le titre « Résumé », et il est composé
% en petit. Ce n'est pas un ornement : c'est le seul passage du document qui
% ne soit pas des mathématiques, et un lecteur doit pouvoir le reconnaître
% avant de l'avoir lu — puis le sauter s'il connaît déjà le sujet, ou s'y
% arrêter s'il ne le connaît pas. Le filet qui le suit marque la reprise du
% corps.

\newenvironment{resume}
  {\small\setlength{\parskip}{0.45em}}
  {\par\medskip\hrule\medskip}

% --- Mots-clés ---------------------------------------------------------------
%
% En anglais, en clôture du résumé de la lecture modernisée : le vocabulaire
% moderne sous lequel chercher ce que les feuillets construisent. Le manifeste
% du site (`npm run manifest`) les extrait pour étiqueter la cote entière —
% cette ligne est la source unique des tags, il n'y a pas de fichier à côté.
\newcommand{\keywords}[1]{\par\smallskip\noindent
  {\footnotesize\textsc{Keywords} --- \textit{#1}}\par}

% --- L'appareil critique ----------------------------------------------------

% Le numéro de feuillet, en marge. C'est l'ancre qui permet de régler un
% désaccord sur une formule en regardant *un* feuillet plutôt qu'une chemise
% de six cents. Le site s'en sert aussi pour tourner les pages du fac-similé
% quand on fait défiler la transcription.
\newcommand{\leaf}[1]{%
  \marginnote{\footnotesize\color{gray!70}\textsf{#1}}%
  \label{leaf:#1}\ignorespaces}

% Illisible : on ne devine pas. Un mot inventé qui se lit comme les autres est
% le pire résultat possible.
\newcommand{\ill}{\textcolor{red!60!black}{[\,\ldots\,]}}

% Lecture douteuse : le mot est donné, et signalé comme incertain.
\newcommand{\uncertain}[1]{\uline{#1}}

% Ajout de l'éditeur — une lettre manquante, un mot sous-entendu.
\newcommand{\add}[1]{\textcolor{blue!50!black}{[#1]}}

% Biffé par Grothendieck lui-même. Ce qu'il a rayé fait partie du document :
% une rature dit souvent où la pensée a changé de direction.
\newcommand{\struck}[1]{\sout{#1}}

% Note du transcripteur, sur l'état matériel du feuillet.
\newcommand{\note}[1]{\par\smallskip{\small\color{gray!60!black}\textit{[#1]}}\par\smallskip}

% Note marginale de Grothendieck, distincte du corps du texte.
\newcommand{\marginal}[1]{\par\smallskip\noindent\hspace{1em}%
  {\small\color{gray!60!black}#1}\par\smallskip}

% --- Titre ------------------------------------------------------------------

\AtBeginDocument{%
  \begin{center}
    {\large\bfseries Fonds Alexandre Grothendieck --- cote n\textsuperscript{o}~\grothFolder}\\[2pt]
    {\itshape\grothFoldertitle}\\[4pt]
    {\small feuillets \grothFirst--\grothLast\ (lot \grothBatch)}\\[2pt]
    {\footnotesize\color{gray!60!black}Transcription établie d'après le fac-similé numérisé
     par l'Université de Montpellier.\\
     Les lectures incertaines sont soulignées, les illisibles notées [\,\ldots\,],
     les ajouts entre crochets.}
  \end{center}
  \vspace{1em}}

\endinput


\folder{162-5}
\batch{2}
\pages{21}{40}
\dating{1968-[à partir de 1970] — le groupe « Dossiers rassemblés par Grothendieck sur différentes thématiques » (161-1 à 162-6) est daté [après 1961-vers 1977]}
\watermark{Édition de démonstration}
\foldertitle{Tapis de Quillen : tapuscrit et copies de tapuscrits (1968, s.d.), notes manuscrites (1968), tiré à part (1968).}

\begin{document}

\page{21}

\section*{tapis Quillen (Notes du 6.9.68)}

\note{ce titre est le sien, écrit et souligné en haut à droite de la page ;
il ouvre dix pages manuscrites numérotées 1 à 11 par paragraphes, sans
pagination de sa main}

\subsection*{1. Catégories formelles.}

Catégorie formelle \add{$\mathcal{C}$} « sur $X_0$ » i.e. dont
l'\uncertain{objet} des objets est le schéma $X_0$ (regardé comme
\struck{\ill{}} schéma formel i.e. \emph{ind-schéma} particulier).

\marginal{$X_2$ (trois flèches descendantes, deux montantes) $X_1$ (deux
descendantes, une montante) $X_0$}
\note{en marge gauche, le schéma d'un objet semi-simplicial tronqué
$X_2 \to X_1 \rightrightarrows X_0$, faces et dégénérescences, dessiné
verticalement}

Revient à la donnée d'un pro-Anneau strict \add{« admissible »} $A$ sur
$X_0$, augmenté vers $\mathcal{O}_{X_0}$, l'idéal d'augmentation étant idéal de
définition, avec deux structures \add{de bi-$\mathcal{O}_{X_0}$-proanneaux} de
$\mathcal{O}_{X_0}$-modules (g. et dr.), et d'un homom.
\[
A \xrightarrow{\ \Delta\ } A_{\mathrm{dr}} \otimes_{\mathcal{O}_{X_0}} {}_{\mathrm{g}}A
\]
\struck{\ill{}}, satisfaisant des axiomes évidents a) « associativité » de
$\Delta$ b) unitarité g. et dr. vis-à-vis de l'augmentation. [Cela
serait-il nécessairement un \struck{catégorie} \emph{groupoïde} formel, i.e.
y a-t-il une \emph{symétrie} $s$ de $A$, jouant vis-à-vis de $\Delta$ et de
l'augmentation $\varepsilon$ le rôle \uncertain{de passage à l'inverse} ?]

On voit $\Omega_{\mathcal{C}}$ = faisceau conormal de $X_0$ dans $X_1$ ;
c'\uncertain{est} \uncertain{a priori} un pro-faisceau qu. coh. sur $X_0$,
mais dans le cas où $A$ est $I$-\emph{adique} [$I$ l'idéal
d'augmentation] c'est un module qu. coh. ordinaire.

On définit dans $\overset{*}{\Lambda}\Omega$ une différentielle d'anneau
gradué \add{$d = d_{\mathcal{C}}$}. Sa valeur sur
$\overset{0}{\Lambda}\Omega = \mathcal{O}_{X_0} \to \Omega$ :
\[
df = \mathrm{pr}_2^*(f) - \mathrm{pr}_1^*(f) \mod I^2 .
\]

\marginal{\uncertain{infinitésimale} \uncertain{vu plus bas}}
\emph{Sorite} : \uncertain{développer} : $A = (A_\alpha)$, regardons $A_\alpha$ comme
module par la structure gauche, alors l'application
$f \mapsto \mathrm{pr}_2^*(f) : \mathcal{O}_X \to A_\alpha$ est un op. diff.
[d'ordre $n$ si $I_\alpha^{n+1} = 0$], et par suite \uncertain{les}
\struck{\ill{}} sections de
$\underline{\mathrm{Hom}}_{\mathcal{O}_X}(A_\alpha, \mathcal{O}_X) =
\mathcal{D}^{\mathcal{C}}_\alpha$ définissent un op. différentiel dans
$\mathcal{O}_X$. Alors
$\mathcal{D}^{\mathcal{C}} = \varinjlim_\alpha \mathcal{D}^{\mathcal{C}}_\alpha$
est un faisceau d'algèbres \add{unitaires ass.} [grâce à l'application
$\Delta$], qui s'envoie dans le faisceau des op. différentiels [qui
\uncertain{correspond}

\page{22}
au fond ici ce serait la formalisation de la catégorie \ill{} définie par
$X_0/\mathrm{Spec}\,\mathbb{Z}$). En particulier
\[
\mathfrak{g}^{\mathcal{C}} = \struck{\mathrm{d\acute{e}f}}\;
\underline{\mathrm{Hom}}_{\mathcal{O}_{X_0}}(\Omega_{\mathcal{C}}, \mathcal{O}_X)
= \varinjlim_{\alpha} \underline{\mathrm{Hom}}_{\mathcal{O}_{X_0}}(\Omega_{\mathcal{C}_\alpha}, \mathcal{O}_{X_0})
\]
opère sur $\mathcal{O}_{X_0}$ par \emph{dérivations} …. Cf aussi plus bas
pour \uncertain{sa} dérivation à partir des complexes de De Rham
$\Omega^{*}_{\mathcal{C}}$ de $\mathcal{C}$.
\note{la lettre gothique $\mathfrak{g}$ est celle de la page ; elle désigne
ici, comme aux pages 23 et suivantes, le faisceau
$\underline{\mathrm{Hom}}(\Omega, \mathcal{O})$ des champs de vecteurs de la
catégorie formelle}

\subsection*{2. Complexes de De Rham généraux.}

Appelons complexe de De Rham \uncertain{abstrait} sur un \struck{\ill{}}
Module \add{\uncertain{q. coh.}} $\Omega$ sur un schéma $X_0$ [NB
\uncertain{considérer} aussi \ill{} espaces annelés, \ill{}
\add{tout complexe obtenu par} les \uncertain{sorites} de 1.] une
\struck{\ill{}} différentielle d'anneau gradué sur
$\overset{*}{\Lambda}\Omega$.

\emph{Remarque 1}. Cette différentielle est connue quand on la connaît en
degrés 0 et 1 :
\[
\delta_0 : \mathcal{O}_{X_0} \longrightarrow \Omega \qquad
\text{[i.e. } \Omega^1_{X_0/\mathbb{Z}} \xrightarrow{\ u\ } \Omega
\text{ par } \delta(f) = u(df)\text{ ]}
\]
\add{une dérivation}, i.e.
\[
\alpha)\quad \delta_0(fg) = \delta_0(f)\,g + f\,\delta_0(g)
\]
\[
\delta_1 : \Omega \longrightarrow \overset{2}{\Lambda}\Omega
\]
Conditions nécessaires sur $\delta_1$ :
\[
\begin{cases}
\beta)\quad \delta_1(f\omega) = \delta_0(f)\,\delta_1(\omega) + f\,\delta_1(\omega) \\
\gamma)\quad \delta_1\delta_0 = 0
\end{cases}
\]
\note{la page écrit bien $\delta_0(f)\,\delta_1(\omega)$ au premier terme de
$\beta)$, là où la règle de Leibniz donne $\delta_0(f) \wedge \omega$ ; on
transcrit tel quel}

De plus, on peut choisir $\delta_0$ et $\delta_1$ arbitrairement, de façon à
satisfaire ces conditions : ils se prolongent \uncertain{(de façon}

\page{23}
\marginal{\uncertain{dégager un lemme circonstancié} ….}
\uncertain{uniques)} en une \struck{\ill{}} \supplied{anti}dérivation d'Algèbre
\add{\uncertain{grâce} \struck{\uncertain{à la première}} \uncertain{aux}
conditions \uncertain{$\alpha)$ $\beta)$}}, et cette dernière est de
carré nul \uncertain{grâce à} la condition $\gamma)$.

\subsection*{3. \struck{Algèbre de} Anneau de Lie associé à un complexe de De Rham}

Soit $\mathfrak{g} = \underline{\mathrm{Hom}}(\Omega, \mathcal{O}_{X_0})$,
c'est un ind-module quasi-coh. sur $X_0$, qui s'identifie à un Module
qu.-coh. \uncertain{bien} défini, si $X, Y$ sont des sections de
$\mathfrak{g}$ \struck{\ill{}} sur un $U$
\[
[X,Y]\ \struck{(U)} \in \Gamma(U, \mathfrak{g})
\]
avec
\[
(*)\qquad \langle \omega, [X,Y] \rangle
= -\langle \delta\omega, X \wedge Y \rangle + \theta_X \langle \omega, Y \rangle
- \theta_Y \langle \omega, X \rangle ,
\qquad \omega \in \Gamma(U, \Omega)
\]
\note{les deux signes qui précèdent $\theta_X$ et $\theta_Y$ sont récrits
à l'encre par-dessus d'autres ; on lit $+$ puis $-$. Un mot biffé, illisible,
suit $[X,Y]$ au premier membre}

où l'action $X \mapsto \theta_X$ de $\mathfrak{g}$ sur $\mathcal{O}_X$
\add{(par dérivations)} est définie par
\[
\theta_X(f) = \langle \delta f, X \rangle ,
\qquad
\text{[NB } \boxed{\theta_{gX} = g\,\theta_X} \text{ ]}
\]
(résulte de $\alpha)$ : $\theta_X(fg) = \theta_X(f)\,g + f\,\theta_X(g)$)

On constate \add{grâce à $\alpha)$ et $\beta)$} que le deuxième membre est
bien \emph{$\mathcal{O}_{X_0}$-linéaire} en $\omega$, donc définit une
section $X \wedge Y$.
\note{« section $X \wedge Y$ » est tel quel ; on attend une section de
$\mathfrak{g}$, qui est $[X,Y]$}

De plus, on a, utilisant $\gamma)$ :
\[
(**)\qquad \boxed{\theta_{[X,Y]} = [\theta_X, \theta_Y]}
\;\overset{\mathrm{df}}{=}\; \theta_X\theta_Y - \theta_Y\theta_X
\]
\struck{[\ill{} l'autre définition]}

[Sans $\gamma)$, le \struck{\ill{}} \add{\ill{}} différerait au
\uncertain{deuxième membre} \ill{}
$f \mapsto -\langle \delta\delta f, X \wedge Y \rangle$ ; donc si $\Omega$
loc. libre, la condition $\delta\delta = 0$ est \emph{nécessaire} pour
qu'on ait $(**)$]. On vérifie aussi \add{loi de Lie} :
\[
\boxed{[X,X] = 0}\ \text{(\uncertain{trivial})},
\qquad
\boxed{[[X,Y],Z] + [[Y,Z],X] + [[Z,X],Y] = 0}
\]

\page{24}
\[
\boxed{[X, fY] = \theta_X(f)\,Y + f\,[X,Y]}
\qquad \text{résulte de } \alpha)
\]

Prouvons enfin la formule \ill{}. Soit $\omega \in \Gamma(\Omega)$,
\uncertain{d'où} \struck{\ill{}} calculs
\begin{align*}
&\sum_{\substack{\text{perm. circ.}\\ \text{de } X,Y,Z}}
\Big[ -\langle \delta\omega, [X,Y] \wedge Z \rangle
+ \theta_{[X,Y]} \langle \omega, Z \rangle
- \theta_Z \langle \omega, [X,Y] \rangle \Big] \\
&= \sum \Big[ -\langle i_Z \delta\omega, [X,Y] \rangle
+ (\theta_X\theta_Y - \theta_Y\theta_X) \langle \omega, Z \rangle \\
&\qquad + \theta_Z \langle \delta\omega, X \wedge Y \rangle
- \theta_Z\theta_X \langle \omega, Y \rangle
+ \theta_Z\theta_Y \langle \omega, X \rangle \Big] \\
&= \sum -\langle i_Z \delta\omega, [X,Y] \rangle
+ \theta_Z \langle \delta\omega, X \wedge Y \rangle \\
&= \sum \langle \delta i_Z \delta\omega, X \wedge Y \rangle
- \theta_X i_Y i_Z \delta\omega + \theta_Y i_X i_Z \delta\omega
+ \theta_Z \langle \delta\omega, X \wedge Y \rangle \\
&= \sum \big( i_X i_Y \delta i_Z + \theta_Z i_X i_Y
- \theta_X i_Y i_Z + \theta_Y i_X i_Z \big) . (\delta\omega)
\end{align*}
\note{calcul très corrigé : plusieurs termes y sont biffés et récrits, qu'on
ne relève pas un à un ; on donne l'état final des lignes. Sous la dernière,
$\delta i_Z$ est entouré}

[On a utilisé $\boxed{\theta_X = i_X \delta + \delta i_X}$ et les formules
analogues en $Y, Z$, qui \uncertain{donnent}
\begin{align*}
&\sum i_Y i_X \delta i_Z + i_Z \delta i_Y i_X + \struck{\ill{}}\, \delta i_Z i_Y i_X
= i_X \delta i_Y i_Z - \delta i_X i_Y i_Z + i_Y \delta i_X i_Z + \delta i_Y i_X i_Z \\
&\sum i_Y i_X \delta i_Z + \big( i_Z \delta i_Y i_X + i_Y \delta i_X i_Z - i_X \delta i_Y i_Z \big)
= 3\, \delta i_X i_Y i_Z \\
&= \struck{9\, \delta i_X i_Y i_Z +}
\end{align*}
On constate similairement
\[
\sum -\theta_X i_Y i_Z + \theta_Y i_X i_Z = -2 \sum \theta_Z i_X i_Y
\]
\struck{$\sum \theta_Z i_X i_Y = -\sum \theta_X i_Y i_Z$}]
\note{tout ce calcul entre crochets est barré de trois longues diagonales}

donc on trouve
\[
(\text{Jacobi})(\omega) = T(\delta\omega), \quad \text{avec} \quad
T = \sum i_X i_Y \delta i_Z - \theta_Z i_X i_Y \;\struck{= \ill{}}
\]
\add{($\theta_Z = i_Z \delta$ \ill{})}

\page{25}
Utilisons
\[
\boxed{\theta_X = i_X \delta + \delta i_X}
\]
\add{(en degrés quelconques), cf plus bas}

il vient alors, \uncertain{on déduit}
\[
T = \sum i_X i_Y \delta i_Z - i_Z \delta i_X i_Y \;\struck{\ill{}}
\]
ou
\[
T(\delta\omega) = \sum [\, i_{X \wedge Y}, \theta_Z \,](\delta\omega)
\]
pourquoi est-ce nul ??

On a donc une \emph{$\mathcal{O}$-Algèbre de Lie} sur $X_0$, i.e. un
\emph{Anneau de Lie} \struck{sur} $\mathfrak{g}$ sur \struck{$\mathcal{O}_{X_0}$}
$X_0$, muni d'un \emph{homom. d'Anneaux} de Lie
\add{$\theta : X \mapsto \theta_X$} dans $\mathfrak{g}_{X_0/\mathbb{Z}}$, et
d'une structure de $\mathcal{O}_X$-Module rendant $\theta$
$\mathcal{O}_{X_0}$-linéaire et satisfaisant la formule
$[X, fY] = \struck{\ill{}}\, f[X,Y] + \theta_X(f)\,Y$.

\subsection*{4. Complexe des Cochaînes associé à une $\mathcal{O}$-algèbre de Lie \add{$\mathfrak{g}$}}

Cf. une note Bourbaki ; on \uncertain{pose} \struck{(\uncertain{supposant}
\ill{} \ill{})}
\[
C^p = \underline{\mathrm{Hom}}(\overset{p}{\Lambda}\mathfrak{g}, \mathcal{O}_{X_0})
\]
et on y définit une différentielle par une formule connue ….

Lorsque $\mathfrak{g}$ provient d'un complexe de De Rham
$(\overset{*}{\Lambda}\Omega, \delta)$, alors l'homom. d'alg. graduées
évident $\overset{*}{\Lambda}\Omega \to C^{*}(\mathfrak{g})$ est
\emph{compatible} avec les \struck{structures} différentielles.

\page{26}
\subsection*{5.}

Les $\mathcal{O}$-Algèbres loc. libres \add{de t.f.} \add{$\mathfrak{g}$}
sur $X_0$ correspondent exactement aux complexes de De Rham sur $X_0$
construits sur des $\Omega$ loc. libres de t.f., par les correspondances de
3° et 4°.

\subsection*{6.}

Revenons au cas d'une catégorie formelle sur $X_0$, définie par
$A = (A_\alpha)$ sur $X_0$ …, et le complexe de De Rham associé
$(\Omega^* = \overset{*}{\Lambda}\Omega, \delta)$.

\emph{Th. de Quillen} \add{\uncertain{lorsque} $X_0$ est car. nulle}. Le
foncteur $A \mapsto (\Omega^*, \delta)$ induit une \emph{équivalence} entre
la catégorie des catégories formelles \add{$I$-adiques} sur $X_0$ dont le
$\Omega$ est \emph{plat}, et la catégorie des \struck{Algèbres}
\add{complexes} de De Rham $(\overset{*}{\Lambda}\Omega, \delta)$ sur $X_0$
dont le $\Omega$ est \emph{plat}. [Les $A$ en question sont ceux pour
lesquels $\mathrm{Sym}\,\Omega \to \mathrm{Gr}(A)$ est un
\uncertain{isomorphisme}, ce qui \struck{\ill{}} \uncertain{signifie} : une
propriété \uncertain{de} « lissité différentielle ».]

Ceci contient le th. de Cartan sur les groupes formels en car. nulle.

Il y a une variante « à puiss. divisées », quand $A = (A_n)_{n \in \mathbb{N}}$
\add{\uncertain{\ill{}}} est muni sur l'idéal d'augmentation d'une structure
\struck{d'algèbre} \uncertain{de} puiss. divisées, qui passe au quotient sur
les $A_n$, et \uncertain{chaque} $I_n$ étant nilpotent d'ordre $n$, et la
filtration étant celle provenant des puiss. divisées. [À vrai dire, dans
cet énoncé il faut \uncertain{toujours} des \uncertain{relations} entre la
structure de puiss. divisées et l'application $\Delta$ …]

\page{27}
\subsection*{7. Coefficients}

\add{\uncertain{sur} catégorie formelle} Soit $M$ un
$\mathcal{O}_{X_0}$-module. La notion de « donnée de descente » sur $M$
relativement à $\mathcal{C}$ [ou « $\mathcal{C}$-stratification », ou
« opération de $\mathcal{C}$ sur $M$ »] est claire. On constate alors
qu'une telle structure permet de mettre sur
\[
\overset{*}{\Lambda}\Omega \otimes_{\mathcal{O}_{X_0}} M
\]
un opérateur différentiel $\delta_M$, tel que
\[
(*)\qquad \delta_M(fx) = \delta(f)\,x \;\struck{\ill{}}\; + (-1)^{\deg f} f\,\delta_M(x) .
\]
[Il suffit d'avoir une donnée de recollement d'ordre $1$]
Un tel opérateur est connu quand on connaît
\[
M \xrightarrow{\ \delta^1_M\ } \Omega \otimes M
\]
satisfaisant $(*)$ \struck{\ill{}} pour $f \in \Gamma\mathcal{O}_{X_0}$
\add{(donc $\deg f = 0$)} \add{\uncertain{sur}
$\overset{*}{\Lambda}\Omega \otimes M$ cf plus bas}. La \uncertain{fait}
que l'opérateur \uncertain{\ill{}} \add{soit de} \uncertain{carré nul} \ill{}
provient de la condition de descente [\uncertain{il suffit qu'elle soit
vraie} : l'ordre 2 …??].
\note{le crochet ouvert après « La » et l'ajout interlinéaire « soit de »
se lisent mal ; la phrase dit en substance que $\delta_M^2 = 0$ résulte
de la condition de descente, avec la question de savoir si l'ordre 2
suffit}

\subsection*{8. Coefficients relatifs à une algèbre de De Rham}

Plus généralement, si $(\overset{*}{\Lambda}\Omega, \delta)$ est une
alg. de De Rham sur $X_0$, on peut regarder les modules $M$ avec une
\emph{connexion} relativement à elle, i.e. un morphisme
\[
M \xrightarrow{\ \delta_M\ } \Omega \otimes M
\]
satisfaisant
\[
\delta_M(fx) = \delta(f) \otimes x + f\,\delta_M(x) .
\]

\page{28}
On peut alors de façon unique prolonger en une \struck{dérivation}
opération $\delta_M$ sur $\overset{*}{\Lambda}\Omega \otimes M$,
satisfaisant
\[
\delta_M(\omega x) = \delta(\omega)\, x + \omega\,(\delta_M x)
\qquad \big(\omega \in \Gamma\overset{*}{\Lambda}\Omega,;
x \in \Gamma\,\overset{*}{\Lambda}\Omega \otimes M\big)
\]
\note{le signe $(-1)^{\deg\omega}$ attendu devant le second terme manque
sur la page}

Pour qu'on ait $\delta_M\delta_M = 0$, il faut et il suffit que ce soit
\struck{\ill{}} en degré 0, i.e. que le composé
\[
M \xrightarrow{\ \delta^0_M\ } \Omega \otimes M
\xrightarrow{\ \delta^1_M\ } \overset{2}{\Lambda}\Omega \otimes M
\]
soit nul. Ce composé s'appelle la \emph{tenseur courbure}.

La donnée de \add{la connexion} $\delta^0_M$ permet de faire opérer
$\mathfrak{g} = \underline{\mathrm{Hom}}_{\mathcal{O}_X}(\Omega, \mathcal{O}_X)$
dans $M$, grâce à
\[
\theta^M_X(x) = (X \otimes \mathrm{id}_M)(\delta^0_M(x))
\]
On a alors
\[
\begin{cases}
a)\quad \theta^M_X(fx) = \theta_X(f)\,x + f\,\theta^M_X(x) \\
b)\quad \theta^M_{fX}(x) = f\,\theta^M_X(x)
\end{cases}
\]
D'autre part, si la connexion est \add{sans courbure}
\struck{\uncertain{intégrable}}, alors il est vrai que
\[
c)\quad \theta^M_{[X,Y]} = [\theta^M_X, \theta^M_Y] .
\]

\page{29}
\subsection*{9. \struck{\ill{} d'une} Opérations d'une $\mathcal{O}$-alg. de Lie}

Soit $\mathfrak{g}$ une $\mathcal{O}$-Algèbre de Lie sur $X_0$, $M$ un
$\mathcal{O}_{X_0}$-Module. On dit que $M$ a une \emph{$\mathfrak{g}$-connexion}
si on fait « opérer » $\mathfrak{g}$ sur $M$ de façon à satisfaire a) et b)
plus haut. On définit alors un tenseur courbure
\[
K \in C^2(\mathfrak{g}, \underline{\mathrm{End}}_{\mathcal{O}_{X_0}}(M))
\]
en posant
\[
K(X,Y) = \theta^M_{[X,Y]} - [\theta^M_X, \theta^M_Y]
\]
Le tenseur courbure est donc nul ssi on a
\[
\theta^M_{[X,Y]} = [\theta^M_X, \theta^M_Y] ,
\]
alors la connexion est dite \emph{intégrable}. Si $\mathfrak{g}$ provient
de $\Omega$ \add{($\mathfrak{g} = \underline{\mathrm{Hom}}_{\mathcal{O}_{X_0}}(\Omega, \mathcal{O}_{X_0})$)},
et $M$ \struck{provient} muni d'une $\Omega$-connexion $\delta^0_M$, d'où
un tenseur courbure
\[
M \longrightarrow \overset{2}{\Lambda}\Omega \otimes M \longrightarrow C^2(\mathfrak{g}, M)
\]
\add{\uncertain{j'aimerais} \ill{}} ; le composé est noté $K'$, et
\[
K(X,Y)(u) = K'(u)(X,X) .
\]
\marginal{Attention signe !}
\note{la page écrit $K'(u)(X,X)$ ; on attend $(X,Y)$}

\subsection*{10.}

Si $(\overset{*}{\Lambda}\Omega, \delta)$ \add{est défini}, avec $\Omega$
loc. libre de t.f. ($\Longleftrightarrow$ $\mathfrak{g}$
\add{$\mathcal{O}$-alg. de Lie avec $\mathfrak{g}$} loc. libre de t.f.)
alors pour un Module $M$ sur $X_0$,

\page{30}
les deux notions de connexion se correspondent biunivoquement, de même que
les deux notions de \add{tenseur courbure, de} connexion intégrable …, et
\uncertain{idem} sur les complexes
\[
\overset{*}{\Lambda}\Omega \otimes M \quad \text{et} \quad C^{*}(\mathfrak{g}, M) .
\]

\subsection*{11.}

Sous les conditions du th. de Quillen, ($X_0$ de car. nulle,
$(\overset{*}{\Lambda}\Omega, \delta)$ \add{avec $\Omega$ plat}
\struck{\ill{}} $\Longleftrightarrow$ correspondant à une catégorie formelle
$\mathcal{C}$ avec $\Omega$ plat), pour un Module $M$ sur $\mathcal{O}_{X_0}$,
les structures de $\mathcal{C}$-descente sur $M$ correspondent aux
\struck{\ill{}} $\Omega$-connexions sur $M$. \struck{De} plus, la cohomologie de
De Rham de $M$, i.e. $\mathbb{H}^{*}(X, \overset{*}{\Lambda}\Omega \otimes M)$,
est canoniquement isomorphe à la coh. de Hochschild (ou Čech-Alexander) de
$\mathcal{C}$ à coeff. \struck{dans} $M$, plus précisément il y a un
\add{canonique} quasi-isom. de \add{complexes de} \struck{faisceaux}
\uncertain{\ill{}} (dans la catégorie dérivée)
\[
\overset{*}{\Lambda}\Omega \otimes M \simeq C^{*}(\mathcal{C}, M)
\]
\marginal{Relier au th. fondamental du dévissage des cristaux …. Peut-il y
avoir ce raisonnement dans \uncertain{un} topos convenable ?!}
\note{la note de marge, écrite en oblique en regard des lignes sur la
cohomologie de Hochschild, est lue sur une vue redressée ; la dernière
ligne en est incertaine}

\emph{Variante à puissances divisées.}

\page{32}
\marginal{notes du 10.9.68}
\note{à partir d'ici et jusqu'au bas de la page 40, une copie du tapuscrit
« Tapis de Quillen » dont les pages 2 à 10 du dossier donnent déjà le texte
(tapuscrit, p. 1 à 9), reprise ici à l'encre bleue ; le texte tapé est le
même, les interventions à l'encre ne le sont pas toutes. Comme pour les
pages 41 et suivantes de cette même copie, le texte tapé, y compris les mots
frappés dans l'interligne, est donné en clair ; les ajouts à l'encre sont
donnés comme ajouts, les mots tapés biffés à la main comme ratures, les
notes à l'encre détachées du texte à part ; les frappes surchargées et les
x de la machine ne sont pas relevés. Les soulignements tapés sont en
italique. Les exposants frappés de $C_0$, $C_1$, $X_{ij}$ sont récrits en
indices à la main, et l'on suit la correction}

\section*{Tapis de Quillen.}

\subsection*{I \quad Relations entre catégories et ensembles semi-simpliciaux.}

A toute catégorie $C$, on associe un ensemble semi-simplicial $S(C)$,
trouvant ainsi un foncteur pleinement fidèle
\[
S : (\mathrm{Cat}) \longrightarrow (\mathrm{Ssimpl}) .
\]
Les systêmes locaux d'ensemble sur $SC$ correspondent aux fancteurs sur $C$
qui transforment toute flêche en isomorphisme (i.e. qui se factorisent par
le groupoïde associé à $C$). Les $H^i$ sur $SC$ d'un tel systême local
($H^0$ pour ens, $H^1$ pour groupes, $H^i$ quelconques pour groupes
abéliens) s'interprètent en termes des foncteurs
$\varprojlim^{(i)}$ dérivés de $\varprojlim$, ou si on préfère, des $H^i$
(du \emph{topos} $C$. On voit ainsi à quelle condition un foncteur
$C \to C'$ induit un homotopisme $SC \to SC'$ : en vertu du critère
cohomologique de Artin-Mazur, il f et s que pour tout systême de
coéfficients $F'$ sur $C'$, l'homomorphisme naturel
$\varprojlim^{(i)}_{C'} F' \to \varprojlim^{(i)}_{C} F$ soit un
isomorphisme (pour les $i$ pour lesquels cela a un sens).

A $C$ on peut associer le topos $\widetilde{C}$, qui varie de façon
\emph{covariante} avec $C$. (NB le foncteur $C \mapsto \widetilde{C}$ n'a
plus rien de pleinement fidèle, semble-t-il ??) Les systêmes de
coéfficients ensemblistes sur $C$ (déf = les foncteurs
$C^{\circ} \to \mathrm{Ens}$ transformant isomorphismes en
isomorphismes) correspondent aux faisceaux localement constants, i.e. les
objets localement constants de $\widetilde{C}$, définis intrinsèquement en
termes de $\widetilde{C}$. Ainsi, le fait pour un foncteur
$F : C \to C'$ d'induire un homotopisme $S(C) \to S(C')$ ne dépend que du
morphisme de topos $\widetilde{F} : \widetilde{C} \to \widetilde{C}'$
induit, et signifie que pour tout faisceau localement constant $F'$ sur
$C'$ i.e. sur $\widetilde{C}'$, les applications induites
$H^i(\widetilde{C}', F') \to H^i(\widetilde{C}, \widetilde{F}^*(F'))$ sont
des isomorphismes (pour les $i$ pour lesquels cela a un sens).
\note{la parenthèse ouverte avant « du \emph{topos} $C$ » n'est pas refermée
sur la page ; « il f et s » (il faut et il suffit) est tel quel}

\page{33}
On a aussi un foncteur évident
\[
T : (\mathrm{Ssimpl}) \longrightarrow (\mathrm{Cat}) ,
\]
en associant à tout ensemble semi-simplicial $X$ la catégorie
$T(X) = \Delta_{/X}$ des simplexes sur $X$, dont l'ensemble des objets est
la réunion disjointe des $X_n$ …. (c'est une catégorie fibrée sur la
catégorie $\Delta$ des simplexes types, à fibres les catégories discrètes
définies par les $X_n$). Ceci posé, Quillen prouve que pour tout $X$,
$ST(X)$ est isomorphe canoniquement à $X$ dans la catégorie homotopique
construite avec $(\mathrm{Ssimpl})$, et que pour toute $C$, la catégorie
$TS(C)$ est canoniquement « homotopiquement équivalente à $C$ » i.e.
canoniquement isomorphe à $C$ dans la catégorie quotient de
$(\mathrm{Cat})$ obtenue en inversant les foncteurs qui sont des
homotopismes. Ces isomorphismes sont fonctoriels en $X$. Il en résulte
formellement qu'un morphisme $f : X \to Y$ dans $(\mathrm{Ssimpl})$ est un
homotopisme si et seulement si il en est ainsi de
$T(f) : T(X) \to T(Y)$, d'où des foncteurs
$S' : (\mathrm{Cat})' \to (\mathrm{Ssimpl})'$ et
$T' : (\mathrm{Ssimpl})' \to (\mathrm{Cat})'$ entre les catégories
« homotopiques » construites avec $(\mathrm{Cat})$ resp.
$(\mathrm{Ssimpl})$, qui sont quasi-inverses l'un de l'autre.

De plus, Quillen construit un isomorphisme canonique et fonctoriel dans
$(\mathrm{Cat})'$ entre $C$ et la catégorie opposée $C^{\circ}$, ou ce qui
revient au même, un isomorphisme canonique et fonctoriel dans
$(\mathrm{Ssimpl})'$ entre $S(C)$ et $S(C^{\circ})$. La définition est telle
que le foncteur induit sur les systèmes locaux sur $C$ transforme le
foncteur contravariant $F$ sur $C$, transformant toute flèche en flèche
inversible, en le foncteur covariant (i.e. contravariant sur $C^{\circ}$)
ayant mêmes valeurs sur les objets, et obtenu sur les flèches en remplaçant
$F(u)$ par $F(u)^{-1}$ ; en d'autres termes, l'effet de l'homotopisme de
Quillen sur les groupoïdes fondamentaux est l'isomorphisme évident entre
les groupoïdes fondamentaux de $C$ et de $C^{\circ}$, compte tenu que le
deuxième est l'opposé du

\page{34}
premier. Comme application, Quillen obtient une interprétation faisceautique
de la cohomologie d'un ensemble semi-simplicial à coéfficients dans un
système local covariant $F$ (défini
classiquement par le complexe cosimplicial des
$C^n(F) = \coprod_{x \in X_n} F(x)$) : on considère le système local
contravariant défini par $F$, on l'interprète comme un faisceau sur
$T(X)$ i.e. objet de $(\mathrm{Ssimpl})_{/X}$, et on prend sa
cohomologie. — Cependant, quand $F$ est un système de coéfficients
covariant pas nécessairement local, on n'a tjrs pas d'interprétation de
ses groupes de cohomologie classiques en termes faisceautiques ; ni,
lorsque $F$ est contravariant, de son homologie, ou inversement de sa
cohomologie faisceautique en termes classiques.
\marginal{$\prod$ ?}
\note{le $\coprod$ de $C^n(F)$ est entouré à l'encre, et la marge en regard
demande $\prod$ ?}

A propos de la notion de foncteur qui est un homotopisme. Quillen montre
qu'un tel foncteur $F : C \to C'$ induit une équivalence entre la
sous-catégorie triangulée $D^b_{\mathrm{lc}}(C')$ de la catégorie
dérivée bornée de celle des faisceaux abéliens sur $C'$, dont
les faisceaux de cohomologie sont des systèmes locaux, et la catégorie
analogue pour $C$ ; et réciproquement. On peut dans cet énoncé introduire
aussi n'importe quel anneau de base (à condition de le supposer $\neq 0$
dans le cas de la réciproque) ; la partie directe vaut aussi avec un
anneau de coefficients pas néc. constant, mais constant tordu. Je pense que
ce résultat (facile) doit pouvoir se généraliser ainsi : Soit
$f : X \to X'$ un morphisme de topos qui soit tel que pour tout faisceau
loc. const. sur $X'$, $f$ induise un isomorphisme sur les cohomologies
(avec cas non commutatif inclus). Supposons que \struck{\ill{}}
\add{$X$ et $X'$} soit \emph{localement homotopiquement trivial}, i.e. que
pour tout entier $n$ \add{$\geq 1$}, tout objet $U$ \struck{de $X'$} ait un
recouvrement par des $U_i \to U$, tels que a) tout système local sur $U$
devient constant sur $U_i$, et toute section sur $U$ devient constante sur
$U_i$, et b) pour tout groupe abélien $G$, les
$H^j(U,G) \to H^j(U_i,G)$ sont nuls pour $1 \leq j \leq n$. Alors le
foncteur $D^b_{\mathrm{lc}}(X') \to D^b_{\mathrm{lc}}(X)$ induit par $f$
est
\marginal{Attention, cette condition n'est \uncertain{typiquement}
\emph{pas} satisfaite par les schémas avec leur topologie étale (mais bien
par \struck{\ill{}} \uncertain{variétés} loc. noeth. avec top. Zariski).}
\note{la note de marge est renvoyée par un signe au bas de la page, en
regard de la définition « localement homotopiquement trivial », que deux
traits verticaux signalent dans la marge gauche}

\page{35}
une équivalence). Même énoncé si on met dans le coup un système local
d'anneaux sur $X'$. Enfin, $f$ induit \uncertain{une équivalence} entre la
catégorie des coefficients locaux sur $C$ et celle des coéff. locaux sur
$C'$. Principe de démonstration : on commence par prouver ce dernier
résultat, en notant que si un topos est localement hom. trivial, il est loc.
connexe et loc. simplement connexe, d'où une bonne théorie du $\pi_0$ et du
$\pi_1$ (qui sont ici discrets), et on est ramené à un cas particulier du
critère d'homotopisme de Artin et Mazur, savoir un critère
\add{cohomologique} pour qu'un homomorhisme de groupes $G \to H$ (ici des
groupes $\pi_1$ de $X, X'$) soit un isomorphisme : il doit induire des
isomorphismes sur les $H^0$ et $H^1$ (y inclus dans le cas non
commutatif…). On prouve la pleine fidélité en se ramenant par la suite
spectrale aux cas des $\mathrm{Ext}^i$ de coéfficients locaux. Par la suite
spectrale encore, cela résultera du fait suivant : si $X$ est localement
hom. trivial, alors la catégorie des faisceaux abéliens loc. constants est
stable par $\underline{\mathrm{Ext}}^i$, et le foncteur $M \mapsto M_X$ de
$(\mathrm{Ab})$ dans $X_{\mathrm{ab}}$ commute aux dits $\mathrm{Ext}^i$.
En fait, $X$ et $X'$ étant loc hom triviaux, les conditions suivantes sur
$f$ seront équivalentes :
\note{le début de la page achève la phrase de la page 34 ; les mots qui
précèdent « une équivalence » y ont été biffés à la machine
(« pleinement fidèle (resp. ») et ceux de la troisième ligne surchargés, de
sorte que la lecture « induit une équivalence » est la plus probable, non
la seule. Les $\pi$ sont repassés à l'encre}

\begin{itemize}
  \item[a)] $f$ est un homotopisme, i.e. induit pour tout système local
  (pas néc commutatif) sur $X'$ un isomorphisme sur les $H^i$.
  \item[b)] $f$ induit une équivalence
  $D^b_{\mathrm{lc}}(X') \to D^b_{\mathrm{lc}}(X)$.
  \item[a')] $f$ induit une équivalence entre la catégorie des systèmes
  locaux abéliens sur $X'$ et $X$, et des isomorphismes sur les $H^i$
  correspondants (donc on ne prend ici que des coéfficients commutatifs).
\end{itemize}

\marginal{?}
J'ignore si on peut dans a) se borner aux systèmes locaux commutatifs.
L'équivalence entre a) et b) fournit une première justification ou
motivation pour définir des types d'homotopie via la catégorie
$D^b_{\mathrm{lc}}(X)$,

\page{36}
éventuellement muni de la sous-catégorie pleine de tous les systèmes
locaux sur $X$, et du foncteur cohomologique sur $D^b_{\mathrm{lc}}(X)$ à
valeurs dans ladite catégorie, et bien sûr du produit tensoriel (mais alors
on sort de $D^b$ pour entrer dans $D^-$, redactor demerdetur).

\subsection*{II \quad $n$-catégories, catégories $n$-uples, et (Gr)-catégories.}

Appelons bi-catégorie une catégorie dans $(\mathrm{Cat})$, i.e. un triple
$(C_0, C_1, p_1, p_2, \pi)$, où $C_0$ est une catégorie,
$p_1, p_2 : C_1 \rightrightarrows C_0$ deux foncteurs,
$\pi : (C_1, p_1) \times_{C_0} (C_1, p_2) \to C_1$ un foncteur, avec les
axiomes habituels. Interprétant $(\mathrm{Cat})$ comme la sous-catégorie
pleine de $(\mathrm{Ssimpl})$ des objets qui …., et de même pour les
catégories à valeurs dans une catégorie telle $(\mathrm{Cat})$, on trouve
que la catégorie $(\mathrm{Bicat})$ des $(\mathrm{Cat})$-catégories
est équivalente à la sous-catégorie pleine de la catégorie des ensembles
bi-semisimpliciaux
\[
\begin{array}{cccc}
X_{00} & X_{01} & \cdots & X_{0n} \\
X_{10} & X_{11} & \cdots & X_{1n} \\
\cdots & & &
\end{array}
\]
formée des objets qui ligne par ligne et colonne par colonne satisfont la
condition d'exctitude familière (transformer sommes amalgamées en produits
fibrés). De même la catégorie $(\mathrm{Tricat})$ des
$(\mathrm{Bicat})$-catégories s'interprète en termes d'objets
tri-semisimpliciaux, etc.

Exemple : soit $C$ une catégorie, posons $C_0 = C$,
$C_1 = \underline{\mathrm{Fl}}(C)$, avec structures $p_1, p_2$, évidentes
(d'où une $(\mathrm{Cat})$-catétorie dont les composantes semi-simpliciales
supérieures sont les catégories $\underline{\mathrm{Fl}}^n(C)$ des
$n$-flêches de $C$ …). L'ensemble bi-semisimplicial correspondant a pour
composantes les ensembles $\mathrm{Hom}(\Delta^n \times \Delta^m, C)$. Par
exemple, les objets de $X_{11}$ sont les carrés commutatifs de flêches de
$C$, qu'on peut interpréter (de deux manières différentes) comme

\page{37}
morphismes de flêches.

Notons une symétrie évidente dans la catégorie $(\mathrm{Bicat})$,
interprétée en termes d'ensemble bi-semisimpliciaux, obtenue en échangeant
les lignes et les colonnes dans de tels ensembles bi… Un ensemble
bisemisimplicial qui … peut donc de deux façons différentes être considéré
comme définissant une bicatégorie. De même le groupe symétrique
$\mathrm{Sym}_n$ opère sur la catégorie $(n\text{-upcat})$ des catégories
$n$-uples ….

Une 2-catégorie peut être définie comme constituée par un ensemble
$C_0$ (l'ensemble des objets), qu'on interprètera comme une catégorie
discrète, plus la donnée pour tout $x, y \in C_0$ d'une catégorie
$\underline{\mathrm{Hom}}(x,y)$, plus \add{des foncteurs} composition entre
celles-ci … Prenant la catégorie somme $C_1$ des
$\underline{\mathrm{Hom}}(x,y)$, On peut exprimer la donnée de ces
foncteurs de composition, et des sources et buts \add{$x, y$} dans les
catégories $\underline{\mathrm{Hom}}(x,y)$, comme définis par des foncteurs
« source » et « but » $C_1 \to C_0$, et un foncteur de composition
$C_1 \times_{C_0} C_1 \to C_1$, les axiomes des 2-catégories correspondant
exactement aux axiomes des $(\mathrm{Cat})$-catégories. De cette façon, la
catégorie des 2-catégories est équivalente à la catégorie des
$(\mathrm{Cat})$-catégories qui ont une « catégorie d'objets » qui est
discrète, ou encore à la catégorie des Bicatégories qui, interprétées comme
ensembles bi-semisimpliciaux, ont comme première ligne un ensemble
semisimplicial discret (i.e. un foncteur constant sur $\Delta$). Il parait
qu'on a une interprétation simple dans le même esprit de la notion de
$n$-catégorie en termes d'ensembles $n$-semi-simpliciaux.

Demander à Quillen relations entre la notion de $m$-Catégorie dans la
catégorie $(n\text{-Cat})$, et celle de ensemble
$(m{+}n)$-semisimplicial.
\note{le paragraphe est signalé par un double trait dans la marge gauche ;
la ligne qui le suit, biffée à la machine, n'est pas relevée}

L'interprétation standard semisimpliciale permet d'interpréter

\page{38}
\marginal{Attention : la catégorie $\underline{\mathrm{Hom}}_{\mathrm{grab}}$,
qui s'identifie à la catégorie définie par un
$R\mathrm{Hom}_{\mathbb{Z}}(C_x, C'_x)$, ce point devant être élucidé !}
\note{note encadrée en tête de page, à l'encre, précédée d'un signe de
renvoi dont on n'a pas retrouvé l'appel sur la page. $C_x$, $C'_x$ ne sont
pas définis ici}

de même les $(\mathrm{Cat})$-groupes (ou groupes à valeurs dans
$(\mathrm{Cat})$) comme étant les groupes semisimpliciaux qui « sont » des
catégories i.e. qui satsifont la condition d'exactitude familière,
permettant d'en déterminer tous les composants à l'aide de $C_0$ et $C_1$.
Appliquant le tapis de Dold, dans le cas commutatif, on trouve que les
$(\mathrm{Cat})$-groupes abéliens forment une catégorie équivalente à celle
des complexes de chaines de groupes abéliens qui sont de longueur 1 :
\[
0 \longleftarrow C_0 \longleftarrow C_1 \longleftarrow 0 .
\]
Quand on explicite la $(\mathrm{Cat})$-groupe abélien associé à un tel
complexe de chaines, on trouve « ma » construction chérie de la « catégorie
de Picard » associée au complexe : objets ceux de $C_0$, flêches de $x$
dans $y$ les $z$ dans $C_1$ tels que $dz = y - x$. On trouve donc à
isomorphisme près tous les groupes abéliens de $(\mathrm{Cat})$ ainsi. Un
complément important est que toute catégorie de Picard
\struck{(resp. et commutative)} est équivalente (au sens technique qu'on ne
précise pas ici) à un vrai $(\mathrm{Cat})$-groupe
\struck{(resp. $(\mathrm{Cat})$-groupe commutatif)}. \add{Quand on
\uncertain{simplifie} mod « \uncertain{Grab} »-équivalence, cela correspond
à travailler dans la catégorie \emph{dérivée} des complexes, seuls
invariants \uncertain{sont} $H_0, H_1$ (\uncertain{ou} $\pi_0, \pi_1$).}
\marginal{\struck{\ill{}} Mais \ill{} si on part d'une « cat. de Picard
commut. », elle ne se \uncertain{réalise} \ill{} que par un
$(\mathrm{Cat})$-g. \uncertain{non} \emph{commut.}}
\note{la marge gauche porte, en regard de ces lignes, un premier bloc de
cinq lignes obliques barré de traits verticaux et resté illisible
(« \uncertain{Qualifier} \ill{} mais cela \uncertain{contredit} les résultats
\ill{} »), puis la note conservée, lue sur une vue redressée. Les deux
parenthèses biffées du texte dactylographié, qui étendaient l'énoncé au cas
commutatif, le sont à l'encre ; la note de marge dit pourquoi}

On notera qu'on a une classification beaucoup plus restreinte que pour les
H-espaces (resp. les H-espaces homotopiquement comm.) mod équivalence
d'homotopie, de façon précise entre les groupes resp. les groupes abéliens
de la catégorie homotopique. La raison en est qu'on est beaucoup plus
restrictif dans le contexte catégorique pour le sens de l'associativité
(resp. et de la commutativité), puisqu'on ne demande pas seulement
l'existence d'isomorphismes d'associativité (qui suffirait pour avoir un
groupe dans la catégorie homotopique
$(\mathrm{Cat})' \simeq (\mathrm{Ssimpl})'$, modulo l'existence d'unités et
d'inverse) mais qu'on suppose \emph{donnés} ces isomorphismes, de façon à
satisfaire à certaines compatibilités : c'est ce qui permet, me
semble-t-il, de prouver qu'à équivalence près, on peut supposer qu'il
s'agit de vrais groupes dans $(\mathrm{Cat})$,
\struck{resp. de vraie commutativité.}

\page{39}
On doit trouver une approximation meilleure d'un cran à la classification
topologique en prenant des $(2\text{-Cat})$-groupes resp. des
$(2\text{-Cat})$-groupes abéliens, ou plus généralement des
« Gr-2-catégories » resp. « Grab-2-catégories ». (NB Une Grab-catégorie
définit une Grab-2catégorie comme une catégorie définit une 2-catégorie ;
du point de vue semi-simplicial, cette opération serait analogue à celle de
« cosquélette tronqué » ; Mais également une Grab-2-catégorie définie en
Grab-catégorie, savoir la catégorie des automorphismes de l'objet unité ;
cette opération serait analogue à celle de « tuage des groupes d'homotopie
en dim $> 1$ »). Pour la notion grossière (Gr ou Grab) l'idée de ma
lettre à Verdier était qu'on doit trouver la même chose que les tronqués
au cran 2, dans la catégorie \emph{homotopique} des H-espaces — i.e.
ce qu'on déduit des H-espaces par tuage des groupes d'homotopie en
dimensions $> 2$.
\note{les deux signes $>$ frappés sont repassés à l'encre, le trait de plume
prolongeant leur branche inférieure ; ce n'est pas un $\geq$ (même lecture
à la photocopie, p.~9, et au dossier 111). Le premier chiffre, « 1 », est
lui aussi retouché}

(Demander à Quillen relations exactes entre objets-groupes et objets
connexes de la catégorie homotopique ponctuée : forment-ils des
catégories équivalentes par le foncteur espace des lacets\supplied{ ?}).
\marginal{non, on a un foncteur \uncertain{pl. fidèle} (?) mais pas
\uncertain{ess.} surjectif ….}
\note{« ponctuée » est tapé en interligne et mis en place, après « homotopique », par un trait à l'encre. Le paragraphe est signalé par un double trait vertical dans la marge ;
la réponse de la marge, lue sur une vue redressée, est d'une autre encre}

Pour les \emph{vrais} groupes abéliens de $(2\text{-Cat})$,
l'approche de Quillen consiste à les identifier d'abord à des groupes
abéliens de la catégorie $(\mathrm{Bicat})$, donc à des groupes abéliens
bi-semisimpliciaux particuliers (conditions sur chaque ligne et chaque
colonne). Appliquant la version bi-semisimpliciale du tapis de Dold, on
trouve que la catégorie de ces objets est équivalente à la catégorie des
bicomplexes de chaines de groupes abéliens qui sont nuls en dehors du
carré $0 \leq i, j \leq 1$. La colonne que la première ligne correspond à
un groupe semisimplicial trivial signifie alors qu'elle n'a que le terme de
degré zéro, i.e. que le bicomplexe est réduit à
\[
\begin{array}{ccccc}
C_{00} & \longleftarrow & 0 & \longleftarrow & 0 \\
\uparrow & & \uparrow & & \\
C_{10} & \longleftarrow & C_{11} & \longleftarrow & 0
\end{array}
\]
\note{« La colonne que la première ligne correspond » est tel quel ; les
flèches du diagramme sont en partie tracées ou repassées à l'encre}

\page{40}
Ces bicomplexes s'identifient aux complexes de chaines de longueur 2. Donc
la catégorie de $(2\text{-Cat})$-groupes abéliens est équivelente à celle
des complexes de chaines de longueur 2. L'explicitation de la 2-catégorie
associée à un tel complexe de chaines est mon interprétation chérie. Il
faudrait encore expliciter 2-Grab-catégorie des homomorphismes (en un sens
à préciser) entre deux telles 2-Grab-catégories strictes, en termes d'un
complexe $R\mathrm{Hom}$.

L'idée qui se dégage de cette approche, et que Quillen on espère va
expliciter, c'est que la notions de groupe abélien dans
$(n\text{-Cat})$ correspond à celle de complexe de chaines de
longueur $n$, en raisonnant soit mod. isomorphisme de part et d'autre, soit
mod. une notion naturelle de $n$-équivalence d'un côté, et dans la
catégorie dérivée de l'autre. Ceci doit pouvoir se préciser en utilisant
des $\mathrm{Hom}^{\bullet}$ resp. des $R\mathrm{Hom}^{\bullet}$.

Si maintenant $C$ est une catégorie avec un « produit tensoriel »
associatif (et éventuellement unitaire), i.e. un bifoncteur avec des
isomorphismes d'associativité (mais sans inverse …), on peut lui associer
sans doute une Gr-catégorie universelle (tout comme à un monoïde on associe
un groupe symmétrisé), d'où des invariants, qu'on pourra désigner par
$K_0(C)$ et $K_1(C)$. On suppose qu'on pourra de même trouver, pour tout
$n$, une Gr-$n$-catégorie universelle où envoyer $C$ de façon compatible
avec $\otimes$ (tout cela dans un sens à préciser), dont les invariants
$\pi_n$ seraient par définition les $K_n(C)$. Un cas intéressant est celui
où $C$ est une catégorie avec produits finis (par exemple une catégorie
additive) et $\otimes$ est le produit. [Mais on notera que dans le cas
d'une catégorie non seulement additive mais abélienne, il y a une notion
plus adéquate du $K_0$, qui devrait correspondre à une notion plus adéquate
de la suite des $K_i$, inspirée du tapis de Roos ; sans doute c'est en
termes de (« vraie ») catégories triangulées
\note{le tapuscrit s'arrête ici au bas de la page 40, au milieu de la
phrase, que la page 41 achève ; le crochet ouvert avant « Mais », tracé sur
une parenthèse tapée, n'est pas refermé ici}

\end{document}
